% interaction/two_party.tex — Role, withRoles, Counterpart, N-party

\section{Two-Party and Multi-Party Interactions}\label{sec:interaction-two-party}

A $\mathsf{Spec}$ describes \emph{what} moves are exchanged but not \emph{who}
makes them.  Two-party structure is introduced as a decoration---a
$\mathsf{Decoration}\;(\lambda\,\_.\;\mathsf{Role})$---rather than as a
separate inductive type.  This means all $\mathsf{Spec}$ infrastructure
(transcripts, append, replicate, state chains) works unchanged under role
annotations.

\subsection{Roles and the \texorpdfstring{$\Sigma/\Pi$}{Sigma/Pi} duality}

\begin{definition}[Role]
  \label{int:role}
  $\mathsf{Role}$ is the two-element type $\{\mathsf{sender},\;
  \mathsf{receiver}\}$.  $\mathsf{Role.swap}$ exchanges the two values;
  it is an involution.
  \lean{Interaction.Role}
\end{definition}

The key observation is that ``choosing a move'' ($\Sigma$-type) and ``responding
to any possible move'' ($\Pi$-type) are dual operations:

\begin{definition}[Action and Dual]
  \label{int:role-action-dual}
  For a move type~$X$, continuation family $\mathit{Cont} : X \to \Type$, and
  monad~$m$:
  \begin{align*}
    \mathsf{Action}\;\mathsf{sender}\;m\;X\;\mathit{Cont}
      &= (x : X) \times m\;(\mathit{Cont}\;x), \\
    \mathsf{Action}\;\mathsf{receiver}\;m\;X\;\mathit{Cont}
      &= (x : X) \to m\;(\mathit{Cont}\;x), \\[4pt]
    \mathsf{Dual}\;\mathsf{sender}\;m\;X\;\mathit{Cont}
      &= (x : X) \to \mathit{Cont}\;x, \\
    \mathsf{Dual}\;\mathsf{receiver}\;m\;X\;\mathit{Cont}
      &= m\;\bigl((x : X) \times \mathit{Cont}\;x\bigr).
  \end{align*}
  \lean{Interaction.Role.Action}
  \lean{Interaction.Role.Dual}
  \uses{int:role}
\end{definition}

\subsection{Role decorations, strategies, and counterparts}

\begin{definition}[RoleDecoration]
  \label{int:role-decoration}
  A \emph{role decoration} is $\mathsf{Decoration}\;(\lambda\,\_.\;
  \mathsf{Role})\;\mathit{spec}$: each node is labeled sender or receiver.
  $\mathsf{swap}$ exchanges all labels (involutive).
  \lean{Interaction.RoleDecoration}
  \uses{int:decoration, int:role}
\end{definition}

\begin{definition}[Strategy.withRoles]
  \label{int:strategy-with-roles}
  The \emph{focal strategy} $\Strategy.\mathsf{withRoles}\;m\;\mathit{spec}\;
  \mathit{roles}\;\mathit{Output}$ applies $\mathsf{Action}$ at each node
  according to its role: the focal party \emph{chooses} at its own nodes and
  \emph{responds} at the other party's nodes.
  \lean{Interaction.Spec.Strategy.withRoles}
  \uses{int:role-action-dual, int:role-decoration}
\end{definition}

\begin{definition}[Counterpart]
  \label{int:counterpart}
  The \emph{counterpart} (or \emph{environment}) $\mathsf{Counterpart}\;m\;
  \mathit{spec}\;\mathit{roles}\;\mathit{Output}$ applies $\mathsf{Dual}$ at
  each node: it observes the focal party's choices and effectfully produces its
  own moves in the monad~$m$.
  \lean{Interaction.Spec.Counterpart}
  \uses{int:role-action-dual, int:role-decoration}
\end{definition}

$\mathsf{runWithRoles}$ executes a focal strategy against a counterpart,
producing the full transcript together with both parties' outputs.

\subsection{Per-node monad variants}

For oracle verifiers, different nodes need different monadic effects: a sender
node (prover message) is observed purely, while a receiver node (verifier
challenge) involves oracle computation.

\begin{definition}[Counterpart.withMonads]
  \label{int:counterpart-with-monads}
  Given a $\mathsf{MonadDecoration}$---a per-node choice of monad---the
  counterpart $\mathsf{Counterpart.withMonads}$ uses the node's own monad at
  each step.
  \lean{Interaction.Spec.Counterpart.withMonads}
  \uses{int:counterpart, int:decoration}
\end{definition}

This is the type that oracle verifiers instantiate (Section~\ref{sec:interaction-oracle}):
sender nodes get $\mathsf{Id}$ (pure observation, since $\mathsf{Id}\;\alpha =
\alpha$ definitionally), and receiver nodes get $\OracleComp$ with accumulated
oracle access.  All generic composition combinators for
$\mathsf{Counterpart.withMonads}$ therefore apply directly to oracle
counterparts.

\subsection{Role-aware refinement}

\begin{definition}[Role.Refine]
  \label{int:role-refine}
  $\mathsf{Role.Refine}\;S\;\mathit{spec}\;\mathit{roles}$ is a
  ``sender-only'' decoration: it carries $S\;X$ at sender nodes and recurses
  directly at receiver nodes, with no padding data.  This avoids the
  $\mathsf{PUnit}$ junk that $\mathsf{Decoration.Over}$ would introduce at
  receiver nodes.
  \lean{Interaction.Role.Refine}
  \uses{int:role-decoration, int:decoration-over}
\end{definition}

$\mathsf{Role.Refine}$ is equivalent to
$\mathsf{Decoration.Over}\;(\lambda\,X\;r.\;\mathsf{SenderData}\;S\;X\;r)$
via an explicit $\mathsf{Equiv}$.  This is the type used for
$\mathsf{OracleDecoration}$ (Section~\ref{sec:interaction-oracle}).

\subsection{Multi-party interactions}

For $N > 2$ parties, we introduce a $\mathsf{PartyDecoration}$---a
$\mathsf{Decoration}\;(\lambda\,\_.\;\mathsf{Party})$ for some type
$\mathsf{Party}$---and project to a $\mathsf{RoleDecoration}$ via a
\emph{resolve} function $\mathsf{Party} \to \mathsf{Role}$ (following the
MPST local type projection pattern).

The formalization includes a three-party knowledge-soundness example (prover,
verifier, extractor) with a comparative evaluation.  The conclusion: the
three-party formulation is more general but not uniformly better.  For
straightline extraction the function-based formulation is simpler; for
simulation and state-restoration extraction, the $N$-party interaction
naturally captures the simulator or extractor as a strategy.
